\section{Question 3}

\subsection{Question 3.1}

\partquestion{a)} The hard-margin SVM primal optimisation problem is given as follows:
\[
\begin{aligned}
\min_{w,b} \quad & \frac{1}{2} ||w||^2\\
\text{subject to} \quad & y_i(w^T x_i + b) \geq 1, \quad i = 1,2,\ldots,n.
\end{aligned}
\]
In terms of the margin, you get the functional margin and the geometric margin. The function margin
is equivalent to:
\[\text{Functional Margin} = y\hat{y} = y(w^Tx_i + b).\]
The functional margin measures if a point is on the correct side of the boundary and how confident the
prediction is. The geometric margin measures the real geometric distance from a point to the hyperplane:
\[\text{Geometric Margin} = \frac{y(w^Tx_i + b)}{||w||}.\]
Hence, the full margin width is given by:
\[\frac{2}{||w||}.\]

\partquestion{b)} The support vectors are the training points that are closest to the hyperplane that separates
the classes. The margin is determined based on these points of which the margin boundaries lie
exactly on these points. Hence, the support vectors determine the position and orientation of the
maximum margin hyperplane that best separates the classes. These points only matter because
they directly affect the margin whereas points far from the margin do not affect the optimal 
separating hyperplane if they are moved or removed.
\newpage
\partquestion{c)} For the soft margin objective, the slack variable measures how much each instance is allowed
to violate the margin. As a result, there are two conflicting objectives in that as we want
to reduce margin violations while ensuring the biggest margin. The $C$ hyperparameter allows
us to define the trade-off between these two objectives. By increasing C, the model 
penalises margin violations more strongly, thus resulting in a smaller margin but less
margin violations. Thus, effectively, it reduces the size of the boundary and hence increases
the risk of overfitting.

\subsection{Question 3.2}


\begin{partquestion2}{a)} The RBF kernel maps the input data into an infinite-dimensional feature space. This is because
the mapping is performed by measuring the similarity between two data points based on their
Euclidean distance. Its mathematical definition is called the Taylor series expansion which allows
the $exp$ term in the kernel to be written as infinite sum of polynomial terms.

The main benefit is that it allows the SVM to learn non-linear decision boundaries observed in the
original input feature space. This is done by transforming the input feature space into a higher
dimensional feature space in which a linear decision boundary can be formed. This is especially
useful for the churn data because it is often very unlikely that features and churn have a 
linear relationship. For example, a customer may only churn if their monthly bill is high \textbf{and}
they have many support calls \textbf{and} they do not have a fixed contract. The RBF kernel
hence allows the SVM to capture these interactions between features without having to manually
create these interactions.
\end{partquestion2}

\begin{partquestion2}{b)}

The \textbf{gamma='scale'} computes the following:

\[\gamma = \frac{1}{d \times Var(X)}\]

where $d$ is the number of features and $Var(X)$ is the variance of the training data.

Since the data was previously standardised, the standard deviation and hence variance of the
data would be 1. As a result $\gamma = \frac{1}{4}$ since there are 4 features.

The gamma parameter defines how far the influence of a single training example reaches. High gamma
values represent 'close' whereas low values represent 'far'. Hence, very large gamma values
result in a decision boundary that will be very tight around individual data point (complex boundary). This leads
to an increased risk of overfitting. A very small gamma value results in smoother and less complex decision boundaries. This
leads to an increased risk of underfitting.
\end{partquestion2}

\begin{partquestion2}{c)}

The code used to execute the grid seach can be seen in the Appendix under Listing \ref{lst:grid}. The \emph{GridSearchCV}
class was used from \emph{sklearn}. In addition, \emph{StratifiedKFold} was also used
to ensure splits have the correct class distributions. Lastly the best parameter configuration
was based on the AUC of the ROC curve.

After running the grid search, the best parameter configuration was found to be:
\[[C = 0.1, \gamma = 1]\]
\end{partquestion2}